% GENERATED FILE. Edit canonical lessons, then run npm run generate:books.
% canonical-source-hash: fnv1a64:95ef703c9d1a447b
% canonical-lessons: HI-C68-country, HI-C68-india, HI-C68-city, HI-C68-language, HI-C68-english

\chapter{Where You Come From}
\label{ch:hi-where-you-come-from}
\hlchaptermodality{75}

\begin{chapteropening}
\textbf{By the end of this chapter:} \emph{I can say what country I am in, what city I live in and what language I speak when I meet somebody.}

Complete HI-C68-english and say all five words in order, then name the two languages you can put beside \hi{भाषा}.
\end{chapteropening}

\section[desh]{\hi{देश} (desh) — a country}

\label{lesson:HI-C68-country}



\begin{quote}
Two lessons back you learnt to ask with nothing but your voice, and one lesson back to hang a check on the end: say \textbf{\hi{चाय} \hi{ठीक} \hi{है}} flat, then with the end lifted, then say \textbf{\hi{चाय} \hi{ठीक} \hi{है} \hi{ना}?} and name what each one does. You have also filled in a \hi{नाम} line without a model, and a form asks for more than a name. The next five words are what it asks for.
\end{quote}

\subsection*{You'll want to know: \hi{देश}}
\textbf{\hi{देश}} (\emph{desh}) — \textquotedblleft{}a country\textquotedblright{}.

Sanskrit \hi{देश} (\emph{deśa}) meant a place, a region, a stretch of land — and it comes from the root \hi{दिश्} (\emph{diś-}), \textquotedblleft{}to point out, to show\textquotedblright{}. A \hi{देश} began as whatever a pointing hand took in.

That root travelled. Latin \emph{dīcere}, \textquotedblleft{}to say\textquotedblright{}, and \emph{index}, the pointing finger, grew from the same ancient root as \hi{दिश्}. The finger you use to show someone where you are from is the finger the word is named after.

Every Hindi form you will ever fill in has a \hi{देश} line on it, and every conversation that starts with a greeting reaches it within a minute.

\begin{grammarlens}[title={What you've built}]
The first of five words for saying where you are from.
\end{grammarlens}

\subsection*{Your turn}
\begin{itemize}
\raggedright
  \item \emph{Recall:} say \emph{ghās}, then read \textbf{\hi{कुआँ}} and say what it means
  \item \emph{Hear:} \emph{desh}, said once slowly and once at speed
  \item \emph{Say it:} \emph{desh}
  \item \emph{Say it:} \emph{yah ek desh hai nā?} — last chapter's tag, on your new word
  \item \emph{Read it:} \hi{देश}, then say what it means
\end{itemize}

\begin{tcolorbox}[breakable,colback=teal!4,colframe=teal!35!black,title={Before you move on}]
What does \hi{देश} mean? (\textquotedblleft{}a country\textquotedblright{}.) Say it once more.
\end{tcolorbox}
\section[Bhārat]{\hi{भारत} (Bhārat) — India}

\label{lesson:HI-C68-india}



\begin{quote}
Say \emph{desh}, then write \textbf{\hi{देश}} from memory. Three letters and one matra: \hi{द}, then \hi{े} on top of it, then \hi{श}. You have been taught all three.
\end{quote}

\subsection*{You'll want to know: \hi{भारत}}
\textbf{\hi{भारत}} (\emph{Bhārat}) — \textquotedblleft{}India\textquotedblright{}, the country's name in its own language.

The name is older than the country. \hi{भरत} (\emph{Bharata}) is a figure the Sanskrit epics name, and \hi{भारत} is the form meaning \textquotedblleft{}belonging to Bharata\textquotedblright{} — the land of his people. English \emph{India} took a different route entirely, through Greek and Persian from \hi{सिंधु} (\emph{Sindhu}), the river.

So one country carries two names that started as two different things: a people and a river. In Hindi you will hear \hi{भारत} on the news and \hi{हिंदुस्तान} in the street, and both mean the same place.

\begin{grammarlens}[title={What you've built}]
A \hi{देश} and the name of one: \hi{भारत} \hi{एक} \hi{देश} \hi{है}.
\end{grammarlens}

\subsection*{Your turn}
\begin{itemize}
\raggedright
  \item \emph{Recall:} read \textbf{\hi{कुआँ}}, then say it without looking
  \item \emph{Say it:} \emph{Bhārat}
  \item \emph{Say it:} \emph{Bhārat ek desh hai}
  \item \emph{Say it:} the same sentence twice — flat, then with the end lifted — and then \emph{Bhārat ek desh hai nā?}
  \item \emph{Write it:} a \hi{नाम} line with no model, then \hi{देश} on the line under it
  \item \emph{Write it:} \hi{देश} once more, without looking
\end{itemize}

\begin{tcolorbox}[breakable,colback=teal!4,colframe=teal!35!black,title={Before you move on}]
What does \hi{भारत} mean? (\textquotedblleft{}India\textquotedblright{}.) And \hi{देश}? (\textquotedblleft{}a country\textquotedblright{}.)
\end{tcolorbox}
\section[shahar]{\hi{शहर} (shahar) — a city, a town}

\label{lesson:HI-C68-city}



\begin{quote}
Which of the two words you have is the name of a country, and which is the word for any country at all?
\end{quote}

\subsection*{You'll want to know: \hi{शहर}}
\textbf{\hi{शहर}} (\emph{shahar}) — \textquotedblleft{}a city, a town\textquotedblright{}.

Persian \ar{شهر} (\emph{shahr}), \textquotedblleft{}city\textquotedblright{}, and behind it Old Persian \emph{xšaça}, \textquotedblleft{}kingdom, dominion\textquotedblright{} — a \hi{शहर} was first the seat of a king rather than a crowd of people. The same word ends dozens of place names across the region.

Set it against \textbf{\hi{गाँव}} (\emph{gā̃v}), the village you already own. \hi{गाँव} and \hi{शहर} are the pair Hindi uses for the whole of settled life, and a speaker will place themselves in one or the other within a sentence of meeting you.

A form asks for your \hi{शहर} on the line under your \hi{देश}.

\begin{grammarlens}[title={What you've built}]
Three: \hi{देश}, \hi{भारत}, \hi{शहर}. A country, its name, and the place in it you live.
\end{grammarlens}

\subsection*{Your turn}
\begin{itemize}
\raggedright
  \item \emph{Hear:} \emph{shahar}, then \emph{gā̃v}
  \item \emph{Say it:} \emph{shahar}
  \item \emph{Contrast:} \emph{shahar} and \emph{gā̃v}, and say which one is bigger
\end{itemize}

\begin{tcolorbox}[breakable,colback=teal!4,colframe=teal!35!black,title={Before you move on}]
What does \hi{शहर} mean? (\textquotedblleft{}a city\textquotedblright{}.) Say it once more.
\end{tcolorbox}
\section[bhāshā]{\hi{भाषा} (bhāshā) — a language}

\label{lesson:HI-C68-language}



\begin{quote}
Write \textbf{\hi{शहर}} from memory: \hi{श}, then \hi{ह}, then \hi{र}. No matras at all in this one — three bare letters in a row.
\end{quote}

\subsection*{You'll want to know: \hi{भाषा}}
\textbf{\hi{भाषा}} (\emph{bhāshā}) — \textquotedblleft{}a language\textquotedblright{}.

Sanskrit \hi{भाषा} (\emph{bhāṣā}), \textquotedblleft{}speech, what is spoken\textquotedblright{}, from \hi{भाष्} (\emph{bhāṣ-}), \textquotedblleft{}to speak\textquotedblright{}. The word names the act before it names the thing: a \hi{भाषा} is speaking, turned into a noun.

You have been saying \textbf{\hi{मैं} \hi{हिंदी} \hi{बोलता} \hi{हूँ}} for a long time. \hi{भाषा} is the word that lets you talk about \hi{हिंदी} rather than only in it — and it is a printed line on both of the writing papers you will one day sit, right under \hi{देश} and \hi{शहर}.

\begin{grammarlens}[title={What you've built}]
Four of the five lines a form asks for: \hi{देश}, \hi{भारत}, \hi{शहर}, \hi{भाषा}.
\end{grammarlens}

\subsection*{Your turn}
\begin{itemize}
\raggedright
  \item \emph{Say it:} \emph{bhāshā}
  \item \emph{Say it:} \emph{hindī ek bhāshā hai}
  \item \emph{Write it:} \hi{शहर} once more, without looking
\end{itemize}

\begin{tcolorbox}[breakable,colback=teal!4,colframe=teal!35!black,title={Before you move on}]
What does \hi{भाषा} mean? (\textquotedblleft{}a language\textquotedblright{}.) Say it once more.
\end{tcolorbox}
\section[angrezī]{\hi{अंग्रेज़ी} (angrezī) — English, the language}

\label{lesson:HI-C68-english}



\begin{quote}
Name the four words of this chapter so far.
\end{quote}

\subsection*{You'll want to know: \hi{अंग्रेज़ी}}
\textbf{\hi{अंग्रेज़ी}} (\emph{angrezī}) — \textquotedblleft{}English\textquotedblright{}, the language.

The route is long and worth following. Portuguese traders said \emph{inglês}; Persian took it as \ar{انگریز} (\emph{angrez}), \textquotedblleft{}an Englishman\textquotedblright{}; Hindi added the ending \textbf{-\hi{ी}} that turns a people into their language, and \hi{अंग्रेज़ी} came out. The nuqta under \hi{ज़} is the \emph{z} you learnt with the letters.

Set it beside \textbf{\hi{हिंदी}}, which is built the same way: \hi{हिंद}, the land, plus the same ending. \hi{भाषा} is the word for the class of thing; \hi{हिंदी} and \hi{अंग्रेज़ी} are two members of it.

Both listening papers you will one day sit key on somebody speaking two of these, so the pair matters more than either word alone.

\begin{grammarlens}[title={What you've built}]
Five: \hi{देश}, \hi{भारत}, \hi{शहर}, \hi{भाषा}, \hi{अंग्रेज़ी}. Enough to say what country you are in, what city you live in, and what you speak.
\end{grammarlens}

\subsection*{Your turn}
\begin{itemize}
\raggedright
  \item \emph{Say it:} \emph{angrezī}
  \item \emph{Say it:} \emph{hindī aur angrezī, do bhāshā}
  \item \emph{Write it:} \hi{भाषा} from memory, then \hi{अंग्रेज़ी} beside the printed model
  \item \emph{Say it:} all five in order
\end{itemize}

\begin{tcolorbox}[breakable,colback=teal!4,colframe=teal!35!black,title={Before you move on}]
What does \hi{अंग्रेज़ी} mean? (\textquotedblleft{}English\textquotedblright{}.) And what is the word for a language itself? (\emph{bhāshā}.)
\end{tcolorbox}
